\chapter{Complete Real-World Project: Task Management System}

Summary of the full application built across preceding chapters --- the Task Management System in its complete, production-ready form.

\section{Features}

\begin{tabularx}{\textwidth}{@{}>{\bfseries}p{3.2cm}X@{}}
\toprule
Area & Capabilities \\
\midrule
Authentication & Register, login, logout, Google OAuth login, get
current user, protected routes, role-based authorization \\
Tasks & Full CRUD, assign to a user, status, priority, due date, search,
filter, sort, pagination \\
User & Profile view, avatar upload, update profile \\
Dashboard & Total / completed / pending / in-progress counts, per-user
statistics \\
\bottomrule
\end{tabularx}

\section{Technical Stack}

\begin{tabularx}{\textwidth}{@{}>{\bfseries}p{3.2cm}X@{}}
\toprule
Layer & Technologies \\
\midrule
Frontend & React, Vite, Tailwind CSS, React Router, Axios, Context API \\
Backend & Node.js, Express, MongoDB, Mongoose \\
Auth \& Security & JWT, bcrypt, HttpOnly cookies, CORS, validation, error
handling \\
Media & Cloudinary (image uploads via Multer) \\
\bottomrule
\end{tabularx}

\section{Complete Project Structure}

\begin{diagrambox}[Task Management System -- Folder Tree]
\begin{verbatim}
frontend/
  src/
    components/     # Button, TaskCard, Modal, Navbar, ...
    pages/          # Login, Register, Dashboard, TaskList, ...
    layouts/        # AppLayout, AuthLayout
    hooks/          # useAuth, useTasks, ...
    context/        # AuthContext
    services/       # business logic wrappers around api/
    api/            # axios instance + endpoint functions
    utils/          # formatDate, validators, ...
  package.json

backend/
  src/
    config/         # db.js, cloudinary.js
    controllers/    # authController.js, taskController.js
    middleware/      # protect.js, authorize.js, errorHandler.js
    models/         # User.js, Task.js
    routes/         # authRoutes.js, taskRoutes.js
    services/       # taskService.js
    validators/     # taskValidator.js
    utils/          # generateToken.js
  app.js
  server.js
  package.json
\end{verbatim}
\end{diagrambox}

\section{Purpose of Each Piece}

\begin{tabularx}{\textwidth}{@{}>{\ttfamily}p{3.2cm}X@{}}
\toprule
Path & Purpose \\
\midrule
components/ & Small reusable UI pieces shared across pages \\
pages/ & Route-level screens composed from components \\
layouts/ & Shared page shells (navbar/sidebar wrapping child routes) \\
hooks/ & Reusable stateful logic extracted out of components \\
context/ & Global state (auth) shared without prop drilling \\
api/ & Raw Axios calls, one function per endpoint \\
services/ (frontend) & Optional layer that composes/transforms api calls
for pages \\
utils/ (frontend) & Pure helper functions (formatting, validation) \\
config/ & Backend startup configuration (DB connection, Cloudinary) \\
controllers/ & Request handlers -- parse input, call services, send
response \\
middleware/ & Cross-cutting request logic (auth, error handling) \\
models/ & Mongoose schemas defining data shape and rules \\
routes/ & Maps HTTP method + path to a controller function \\
services/ (backend) & Business logic and database operations, controller
stays thin \\
validators/ & Input validation rule definitions \\
utils/ (backend) & Small backend helpers (e.g. JWT signing) \\
app.js & Configures Express: middleware, routes, error handler \\
server.js & Connects to MongoDB, then starts the HTTP server \\
\bottomrule
\end{tabularx}

\whatwhyhow
{A single cohesive app --- auth, CRUD, dashboard, file uploads --- built from pieces taught in Chapters 1--45, organized into the layered structures introduced early on.}
{Seeing the whole project assembled makes layer relationships (route $\to$ controller $\to$ service $\to$ model, or page $\to$ hook $\to$ api) concrete.}
{Use this as a template: any new MERN project can start from this folder layout and feature set, then grow from there.}
